\unnumberedchap{Conclusion}
\noindent\rule{\textwidth}{0.4pt}

Pour conclure, nous avons pu voir dans l'état de l'art plusieurs des méthodes utilisées par le passé.
À partir d'archives et de recueil de textes, nous pouvons composer un corpus comparable
dans les langues que nous étudions. Nous pouvons utiliser Wikipedia et les liens inter-langues afin
de créer une base d'articles qui sont des textes comparables. La réalisation d'un outil de web scraping pour
la récupération de textes dans les langues que nous étudions est aussi viable afin de composer notre corpus.
Une fois notre corpus composé, nous pouvons utiliser des méthodes établies dans l'état de l'art afin
d'aligner les corpus et d'identifier des fragments parallèles. Pour ce faire, nous pouvons utiliser des outils
d'alignement basé sur le CLIR, ainsi que des méthodes basées sur le plongement lexical.
Comme par exemple, celle établies par des projets tels que FaDA \citep{lohar_2016} ou STRAND \citep{resnik_2003}. STRAND est
particulièrement intéressant dans notre cas. Ce projet se basant sur le minage de site web. 
Ces outils peuvent être combinés à la réutilisation de segments parallèles présentée par 
\citet{steingrimsson_2023}, afin d'extraire des segments intéressants, utiles, qui auraient d'abord été rejetés.